\chapter{Characterization of Recorded-Output Disagreements}
\label{ch:error}

This chapter further characterizes the recorded-output profiles localized in Chapter~\ref{ch:results}. All analyses use the validation-selected parent-balanced HiED--DA trace on the fixed internal held-out set. A \emph{disagreement} means that a recorded output does not agree with any projected benchmark gold parent under the thesis-wide scoring contract. It does not mean that a psychiatrist has confirmed a clinical diagnostic error.

The ranked differential, criterion states, criterion-compatible set, emitted diagnosis set, and committed primary diagnosis are observable system records. They are not hidden chain-of-thought, and their associations do not establish which model call caused a final disagreement. The purpose of this chapter is therefore descriptive: to show what was present in each retained output view, where benchmark agreement was lost or recovered, and which conclusions the records can and cannot support.

\section{Analysis Population and Evidence Contract}
\label{sec:ch8-analysis-population}

Top-3 candidate detection ($D_3$) and committed-primary agreement ($S$) are available for all 1,000 internal held-out cases. Gold-label inclusion in the criterion-compatible set ($I$) is defined for the 915 checker-eligible cases whose projected benchmark set contains at least one non-\emph{Others} parent represented by a Criterion Checker. The remaining 85 cases are scored as the evaluation-only \emph{Others} category and are not assigned an $I$ value.

The projected benchmark set contains one parent in 912 cases, two parents in 83 cases, and three parents in five cases. All case-level agreement statements therefore follow the any-gold contract: a recorded view agrees when it contains at least one projected benchmark parent. When a directional pair is required, the analysis uses the highest-ranked benchmark parent in the genuine Top-3. This rule gives each case one reproducible pair without treating the other benchmark parents as incorrect.

The frozen audit identifies the exact source trace, verifies its 1,000 records and 915 checker-eligible records, and reproduces the joint $D_3$--$I$--$S$ counts reported in Table~\ref{tab:d3rs-joint}. The trace summaries below use only parent-collapsed codes, counts, and criterion-state labels. No transcript text or checker evidence text is reproduced.

\section{Candidate and Compatibility Disagreement Profiles}
\label{sec:ch8-stagewise-taxonomy}

Table~\ref{tab:d3rs-joint} contains six mutually exclusive profiles among the 915 checker-eligible cases: 34 cases with $(D_3,I,S)=(0,0,0)$, 79 with $(0,1,0)$, 12 with $(1,0,0)$, 12 with $(1,0,1)$, 272 with $(1,1,0)$, and 506 with $(1,1,1)$. These complete profiles are more informative than a single final correctness indicator because they distinguish candidate availability, compatibility inclusion, and committed-primary agreement.

\subsection{Top-3 Candidate Misses}

Across all 1,000 internal held-out cases, 198 contain no benchmark gold parent in the genuine Top-3. Of these, 113 are checker-eligible, whereas the remaining 85 are scored as \emph{Others} and have no corresponding gold-parent checker.

Among the 113 checker-eligible candidate misses, 79 nevertheless contain a benchmark gold parent in the criterion-compatible set. The other 34 contain a benchmark gold parent in neither recorded view. The 79 cases are possible because the Criterion Checkers evaluate the complete configured diagnostic scope rather than only the Diagnostician's Top-3. Compatibility inclusion can therefore recover a benchmark parent that is absent from the ranked differential, although the compatible set does not provide a replacement ranking.

\subsection{Ranked but Absent from the Criterion-Compatible Set}

Twenty-four checker-eligible cases contain a benchmark gold parent in the genuine Top-3 but not in the criterion-compatible set. Under DA, 12 of these cases nevertheless commit a benchmark-consistent primary diagnosis, while the other 12 remain benchmark-discordant at the committed-primary output.

This split establishes an important boundary of the recorded outputs. Criterion-compatible inclusion is not a mandatory gate for the DA committed primary: $(D_3,I,S)=(1,0,1)$ occurs in 12 cases. Conversely, compatibility is not sufficient for benchmark-consistent primary commitment: $(1,1,0)$ occurs in 272 cases. This does not make the criterion record irrelevant. It shows that the compatible set is an unranked study artifact and should not be interpreted as either a necessary or sufficient primary-selection rule.

Because the ranked-but-not-compatible disagreement group contains only 12 cases, its criterion-state composition is treated descriptively in Section~\ref{sec:ch8-criterion-evidence}. The larger 272-case ranked-and-compatible nonselection group is characterized next.

\section{Ranked and Criterion-Compatible but Not Selected}
\label{sec:ch8-primary-selection}
\label{sec:ch8-selection-signatures}

The 272 cases with $(D_3,I,S)=(1,1,0)$ form the largest internal recorded disagreement profile after candidate availability and compatibility inclusion are considered. Figure~\ref{fig:ch8-reference-rank} shows the position of the highest-ranked benchmark gold parent within the genuine three-item differential.

\ThesisFigure[0.72\textwidth]{fig_ch8_reference_rank.pdf}{Rank of the highest-ranked benchmark gold parent among the 272 cases with $D_3=1$, $I=1$, and $S=0$.}{fig:ch8-reference-rank}

The highest-ranked benchmark gold parent appears at rank two in 171 cases (62.9\%) and at rank three in 101 cases (37.1\%). It is therefore the immediate alternative to the leading diagnosis in most cases. Rank proximity does not show why another diagnosis was committed, but it identifies the comparative decision that a future selector would need to resolve.

Table~\ref{tab:ch8-selection-pairs} reports the most frequent directional pairs between the highest-ranked benchmark gold parent and the committed primary diagnosis.

\begin{table}[htbp]
\centering
\caption{Most frequent benchmark-parent to committed-primary pairs in the 272-case ranked-and-compatible nonselection profile.}
\label{tab:ch8-selection-pairs}
\small
\begin{tabularx}{\textwidth}{|>{\raggedright\arraybackslash}X|>{\raggedright\arraybackslash}p{3.2cm}|r|r|}
\hline
Highest-ranked benchmark gold parent & Committed primary & Cases & Share\\
\hline
F41 & F32 & 158 & 58.1\%\\
F39 & F32 & 35 & 12.9\%\\
F32 & F41 & 11 & 4.0\%\\
F41 & F42 & 9 & 3.3\%\\
F51 & F32 & 9 & 3.3\%\\
\hline
\end{tabularx}
\parbox{0.96\textwidth}{\footnotesize
\textit{Note.} Each case contributes one pair based on its highest-ranked benchmark gold parent. Percentages use the 272-case profile as the denominator; only the five most frequent pairs are shown.
}
\end{table}

The five pairs account for 222 of the 272 cases (81.6\%). F41-to-F32 alone accounts for 158 cases (58.1\%), whereas the reverse F32-to-F41 direction occurs in 11 cases (4.0\%). The profile is therefore not evenly distributed across diagnostic directions. It is concentrated toward F32, especially when the highest-ranked benchmark parent is F41.

This concentration is a property of the evaluated synthetic trace and benchmark mapping. It does not establish that F32 is clinically overdiagnosed, that F41 is clinically correct in every listed case, or that one finalizer rule caused the asymmetry. It identifies a high-frequency comparison for targeted review and future clinician adjudication.

\section{Diagnostic-Set Construction Disagreements}
\label{sec:ch8-diagnostic-set}
\label{sec:ch8-set-formation}

Committed-primary agreement and complete diagnostic-set agreement are separate outcomes. Among the 1,000 internal held-out cases, HiED--DA commits a benchmark-consistent primary diagnosis in 518 cases. Of these, 409 also achieve exact equality between the emitted diagnosis set and the projected benchmark set. The remaining 109 cases have primary agreement but diagnostic-set disagreement, as shown in Table~\ref{tab:ch8-diagnostic-set-errors}.

\begin{table}[htbp]
\centering
\caption{Mutually exclusive diagnostic-set disagreement types among 109 cases with benchmark-consistent committed-primary output.}
\label{tab:ch8-diagnostic-set-errors}
\small
\begin{tabularx}{\textwidth}{|>{\raggedright\arraybackslash}X|r|r|}
\hline
\textbf{Diagnostic-set disagreement type} & \textbf{Cases} & \textbf{Share}\\
\hline
Additional emitted diagnosis only & 55 & 50.5\%\\
\hline
Omitted benchmark gold diagnosis only & 52 & 47.7\%\\
\hline
Both additional and omitted diagnoses & 2 & 1.8\%\\
\hline
Total & 109 & 100.0\%\\
\hline
\end{tabularx}
\end{table}

The benchmark is genuinely multilabel in 88 cases: 83 contain two projected parents and five contain three. An omitted benchmark diagnosis is therefore observable when the committed primary agrees with one gold parent but another gold parent is not emitted. An additional emitted diagnosis is a label absent from the benchmark set. Under the benchmark contract this is a set disagreement, but it is not proof that the additional diagnosis is clinically false, unnecessary, or non-comorbid. The benchmark does not adjudicate diagnoses that it does not list.

These results explain why committed Top-1 Accuracy and complete-set measures answer different questions. Exact Match requires equality of the full emitted and benchmark sets. Macro-F1 and Weighted-F1 also depend on which diagnoses are emitted or omitted across cases, not only on the committed primary diagnosis.

\section{Criterion-State Patterns}
\label{sec:ch8-criterion-evidence}

This section examines whether broad checker-derived summaries distinguish the four analytical groups formed from Table~\ref{tab:d3rs-joint}. The summaries are generated system records rather than independent clinical annotations.

\subsection{Compatible-Set Size and Criterion-State Composition}

The criterion-compatible set has a median size of six configured categories, with an interquartile range of five to six, in both the 272-case ranked-and-compatible nonselection group and the 518-case committed-primary-agreement group. Compatible-set size therefore provides no descriptive separation between these groups.

Figure~\ref{fig:ch8-criterion-composition} compares the mean within-case proportions of \texttt{met}, \texttt{not\_met}, and \texttt{insufficient\_\allowbreak{}evidence} states. Each case is first summarized by its own criterion-state proportions, after which group means are calculated so that cases with more criterion entries do not receive greater weight. Corresponding medians and interquartile ranges are reported in Appendix~\ref{app:supporting}, Table~\ref{tab:ch8-criterion-ratios}.

\begin{figure}[htbp]
\centering
\includegraphics[width=1.0\textwidth]{fig_ch8_criterion_state_composition.pdf}
\caption{Mean within-case composition of recorded criterion states across the four mutually exclusive analytical groups among the 915 checker-eligible internal cases.}
\label{fig:ch8-criterion-composition}
\end{figure}

The 272-case nonselection group and the committed-primary-agreement group have closely overlapping compositions. Their mean \texttt{met}/\texttt{not\_met}/\texttt{insufficient\_evidence} proportions are 49.9\%/23.0\%/27.1\% and 49.0\%/23.7\%/27.3\%, respectively. The 113 Top-3 candidate misses are also similar at 48.2\%/22.8\%/28.9\%. The 12-case ranked-but-not-compatible disagreement group has descriptively fewer \texttt{met} states and more unresolved evidence at 38.8\%/28.5\%/32.7\%, but the sample is too small to support a stable group-level conclusion.

The broad summaries therefore do not identify which compatible diagnosis will be committed as primary. This result is consistent with the non-nested design: every configured diagnosis is checked, and the compatible set records thresholded compatibility rather than a comparative primary-diagnosis ranking. Aggregate criterion-state composition also cannot be treated as a causal explanation of hidden model reasoning.

\FloatBarrier
\section{Reproducible Recorded-Output Traces}
\label{sec:ch8-audit-traces}

Table~\ref{tab:audit-trace-examples} presents four deterministic examples from the frozen internal trace. For each exact $D_3$--$I$--$S$ profile, the audit selects the case with the lowest source-line number. This rule replaces unaudited illustrative examples with records that can be located directly by case identifier and JSONL source line. The table uses parent-collapsed output codes and gold-parent criterion-state counts; it contains no transcript text or checker evidence text.

\begingroup
\footnotesize
\setlength{\tabcolsep}{3pt}
\renewcommand{\arraystretch}{1.15}
\begin{longtable}{|>{\raggedright\arraybackslash}p{2.6cm}|>{\raggedright\arraybackslash}p{5.0cm}|>{\raggedright\arraybackslash}p{6.2cm}|}
\caption{Deterministically selected recorded-output traces from four exact internal $D_3$--$I$--$S$ profiles.}
\label{tab:audit-trace-examples}\\
\hline
Profile and source & Recorded diagnostic outputs & Gold-parent criterion summary and interpretation\\
\hline
\endfirsthead
\multicolumn{3}{c}{\tablename~\thetable\ continued}\\
\hline
Profile and source & Recorded diagnostic outputs & Gold-parent criterion summary and interpretation\\
\hline
\endhead
\hline
\multicolumn{3}{r}{Continued on next page}\\
\endfoot
\hline
\endlastfoot

$D_3=0$, $I=0$, $S=0$\newline
Case 302945002\newline
JSONL line 6 &
\textbf{Benchmark:} F31\newline
\textbf{Top-3:} F32, F39, F41\newline
\textbf{Compatible:} F43, F51, F32, F39, F98, F41\newline
\textbf{Committed primary:} F32\newline
\textbf{Emitted set:} F32 &
\textbf{State counts:} 3 \texttt{met}, 5 \texttt{not\_met}, 0 \texttt{insufficient\_evidence}\newline
\textbf{Insufficient-evidence IDs:} none\newline
The benchmark parent is absent from both the ranked Top-3 and compatible set.\\
\hline

$D_3=1$, $I=0$, $S=0$\newline
Case 396572486\newline
JSONL line 26 &
\textbf{Benchmark:} F32\newline
\textbf{Top-3:} F39, F32, F41\newline
\textbf{Compatible:} F51, Z71, F98\newline
\textbf{Committed primary:} F39\newline
\textbf{Emitted set:} F39 &
\textbf{State counts:} 1 \texttt{met}, 6 \texttt{not\_met}, 4 \texttt{insufficient\_evidence}\newline
\textbf{Insufficient-evidence IDs:} B1, C2, C5, C7\newline
The benchmark parent is ranked but absent from the compatible set and is not committed.\\
\hline

$D_3=1$, $I=0$, $S=1$\newline
Case 393127297\newline
JSONL line 46 &
\textbf{Benchmark:} F32\newline
\textbf{Top-3:} F32, F39, F51\newline
\textbf{Compatible:} F51, F98, F39, F41\newline
\textbf{Committed primary:} F32\newline
\textbf{Emitted set:} F32 &
\textbf{State counts:} 0 \texttt{met}, 0 \texttt{not\_met}, 11 \texttt{insufficient\_evidence}\newline
\textbf{Insufficient-evidence IDs:} A, B1--B3, C1--C7\newline
The committed primary agrees with the benchmark although the benchmark parent is absent from the compatible set.\\
\hline

$D_3=1$, $I=1$, $S=0$\newline
Case 342807095\newline
JSONL line 25 &
\textbf{Benchmark:} F32\newline
\textbf{Top-3:} F20, F32, F41\newline
\textbf{Compatible:} F43, F20, F32, F41, F39\newline
\textbf{Committed primary:} F20\newline
\textbf{Emitted set:} F20, F32 &
\textbf{State counts:} 8 \texttt{met}, 1 \texttt{not\_met}, 2 \texttt{insufficient\_evidence}\newline
\textbf{Insufficient-evidence IDs:} C2, C7\newline
F32 is ranked, compatible, and emitted as a non-primary diagnosis, but F20 is committed as primary.\\

\end{longtable}
\renewcommand{\arraystretch}{1.0}
\endgroup

The examples illustrate why the recorded views should not be collapsed into one presumed runtime funnel. The criterion path can retain a diagnosis outside the Top-3, DA can commit a benchmark-consistent primary that is absent from the compatible set, and a benchmark diagnosis may be ranked, compatible, and emitted without becoming the committed primary. These are observable relationships among outputs, not clinical adjudications or explanations of the model's internal reasoning.

\section{Chapter Summary}
\label{sec:ch8-completeness-limitations}

This chapter extended the RQ2 localization analysis from aggregate counts to candidate ranks, diagnostic directions, set construction, criterion-state composition, and reproducible case-level traces. Among the 198 all-case Top-3 misses, 113 are checker-eligible; the criterion path still retains a benchmark parent in 79 of those 113 cases. Among 24 cases in which a benchmark parent is ranked but absent from the compatible set, 12 nevertheless have benchmark-consistent committed-primary output and 12 do not.

The largest internal recorded disagreement profile remains the 272 cases in which a benchmark parent is ranked and compatible but another diagnosis is committed as primary. In this profile, the highest-ranked benchmark parent is second in 171 cases, and F41-to-F32 accounts for 158 directional pairs. The pattern is concentrated and reviewable, but it is not evidence that the finalizer alone caused the disagreement or that the benchmark label is the uniquely correct clinical diagnosis.

Primary selection and diagnostic-set construction are also distinct. Of 518 cases with benchmark-consistent committed-primary output, 109 still disagree with the complete benchmark set. Because the benchmark contains one to three projected parents and does not adjudicate unlisted diagnoses, these set disagreements must remain benchmark-relative rather than being described as confirmed clinical omissions or false comorbidities.

Finally, compatible-set size and broad criterion-state composition do not separate the 272-case profile from committed-primary agreement. Taken together, the results show that HiED's retained outputs can localize different forms of benchmark disagreement, while also showing the limit of that localization: compatibility is neither necessary nor sufficient for benchmark-consistent DA commitment, and the recorded artifacts do not provide a causal account of hidden model reasoning.
